% Generated draft: requires mathematical and source review.
\documentclass[11pt]{article}
\usepackage[margin=25mm]{geometry}
\usepackage{fontspec}
\setmainfont{Latin Modern Roman}
\usepackage{amsmath,amssymb,mathtools,unicode-math}
\setmathfont{Latin Modern Math}
\usepackage{enumitem,needspace,xurl,hyperref}
\setlength{\parindent}{0pt}
\setlength{\parskip}{5pt}
\setlength{\emergencystretch}{3em}
\newcommand{\N}{\mathbb{N}}
\newcommand{\Z}{\mathbb{Z}}
\newcommand{\Q}{\mathbb{Q}}
\newcommand{\R}{\mathbb{R}}
\newcommand{\C}{\mathbb{C}}
\newcommand{\F}{\mathbb{F}}
\newcommand{\A}{\mathbb{A}}
\newcommand{\PP}{\mathbb{P}}
\newcommand{\E}{\mathbb{E}}
\newcommand{\eps}{\varepsilon}
\newcommand{\dd}{\,\mathrm{d}}
\begin{document}
\section*{UnsolvedMath ID: 20001649}
\textbf{Source problem code:} catalogue entry
\subsection{Definitions and mathematical statement}
For a finite symmetric generating set \(S\) of a group \(H\), write \(|h|_S\) for word length, \(B_S(L)=\{h:|h|_S\le L\}\), and
\[
 \omega(H,S)=\lim_{j\to\infty}|B_S(j)|^{1/j}.
\]
Uniform exponential growth of \(H\) means \(\inf_S\omega(H,S)>1\). A uniform bound across a class of subgroups would require one constant greater than one working for all those subgroups and all their finite generating sets. The workshop's phrase ``(uniform) uniform'' does not specify this second quantifier or rank dependence. Its literal request for all two-generated subgroups of \(\operatorname{Out}(F_n)\) needs an exponential-growth qualification: cyclic or finite examples cannot satisfy it. The aid instead gives a conditional certificate in a linearly growing witness regime.

Fix a nonabelian free rank \(r\ge2\) and let the reported Bering exponent be
\[
 N_r=(48r^2-48r+3)\,|\operatorname{GL}(r,\Z/3\Z)|.
\]
Suppose \(H\le\operatorname{Out}(F_r)\), \(S\) is a finite symmetric generating set, \(L\ge1\) is an integer, and there exist \(\alpha,\beta\in B_S(L)\) that are linearly growing and have efficient twist trees \(T_\alpha,T_\beta\) with positive Guirardel intersection:
\[
 i(T_\alpha,T_\beta)>0.
\]
Linear growth means the lengths of iterated conjugacy classes grow at most linearly. Twist trees refer to the Bass--Serre trees of efficient Dehn-twist representatives; Guirardel intersection measures the core in the product of the two trees. Exact representative, power, and efficiency conventions are not supplied in the record, so their applicability must be checked against the named independence theorem.

Under these hypotheses the aid reports
\[
 \omega(H,S)\ge 2^{1/(L N_r)}.
\]
The candidate criterion converts short positive-intersection witnesses into this explicit bound. If, for an intended class of eligible subgroups in fixed rank, every finite symmetric generating set admitted such witnesses with \(L\le L_r\), it would give the common lower bound \(2^{1/(L_rN_r)}>1\) for that class. The aid does not establish this short-witness hypothesis. Since \(N_r\) depends on rank, a bound uniform in fixed rank is not by itself a rank-independent growth constant. Nor is the witness property asserted for every two-generated subgroup or every pair of Dehn twists.

The original text includes Matt Clay's suggested ``baby question'' about bounding a ratio in joint work with Alexandra Pettet, while explicitly failing to recall the ratio. No numerator, denominator, or domain is supplied, so no ratio problem can be formalized uniquely. Martin Bridson suggested linear quotients as a possible way to transfer growth from well-understood images; this is a strategy, not a supplied quotient theorem. The extracted text then runs into a new section, ``Homomorphisms with mapping class groups,'' and Juan Souto's question about non-obvious mapping-class-group embeddings in \(\operatorname{Out}(F_n)\). That adjacent open-ended question is preserved as source context, not treated as part of the quantitative growth criterion; neither its surface parameters nor ``non-obvious'' is defined. The stated growth problem remains reported open in the aid.
\subsection{Short English statement}
The aid bounds growth when a short word ball contains two linearly growing elements with positively intersecting efficient twist trees. A uniform bound on their required word length is still missing. The source also omits the necessary exponential-growth scope and the ratio in its suggested baby question.
\subsection{Sources}
\begin{itemize}
\item Local UnsolvedMath registry, permanent ID \(20001649\), source code \texttt{catalogue entry}; \texttt{aim-geometric-group-theory-notes.json}, \texttt{notes[157]}, tag \texttt{problem}.
\item AIM workshop ``The geometry of the outer automorphism group of a free group,'' Problem 3.7: \url{https://aimath.org/WWN/outerauto/outerauto.pdf}. Source type PDF; statement status \texttt{exact}. The supplied verification attributes the wording to the PDF; only the supplied text was used, and no PDF was accessed or compiled.
\item The aid names Bering's uniform-independence exponent without an authoritative bibliographic citation or complete theorem hypotheses. Source discussion names Matt Clay and Alexandra Pettet for an unspecified ratio, Martin Bridson for the linear-quotient strategy, and Juan Souto for the following mapping-class-group question. No full references are supplied; original literature field is empty.
\item Research attempt 1: \texttt{partial\_result}, status at run \texttt{open}, rigor label \texttt{proved}, novelty confidence low; registry status \texttt{partially\_solved}, classification \texttt{PARTIAL-PROGRESS}. Created and updated 14 August 2026. Machine-generated aid, not peer reviewed; independently verify before citation.
\item Ulam AI / UnsolvedMath Contributors; CC BY 4.0.
\end{itemize}
\end{document}
